\chapter{Research Context}

With the widespread adoption of large language models (LLMs), whose knowledge is limited to the parametric memory encoded during training, Retrieval-Augmented Generation (RAG)~\cite{lewis2021retrievalaugmentedgenerationknowledgeintensivenlp} has emerged as the predominant approach in enriching model outputs with non-parametric memory from external  sources. Vector RAG encodes documents as chunks in the form of embeddings and retrieves the top k documents via semantic similarity given a query and adds them as context to the LLM for generation. 

While vector RAG achieved state of the art performance in query focused summarization (QFS), Edge et al.~\cite{edge2025localglobalgraphrag} highlighted that it falls short when it comes to global or sense making queries like 'what are the main themes in this dataset?'. They proposed GraphRag as a solution to this problem. GraphRag adopts a graph-based approach to question answering. An LLM is first used to extract entities and relationships from source documents and construct a knowledge graph. Communities are then identified within the graph and summarised, enabling the system to answer global queries. However, GraphRAG has its own limitations. Graph construction introduces an expensive upfront cost, and for simple query-focused summarisation, vector RAG is far more efficient.

This highlights the need for adopting hybrid approaches to extending an LLM's knowledge. LinkedIn verfied this by deploying a hybrid model to production for six months. This hybrid approach was shown to reduce the median per-issue resolution time by 28.6\% ~\cite{Xu_2024}. 

Cover adaptoive rah here as a sort of relative work

my novel idea/ contribution.

